\documentclass{amsart}
% For dealing with references we use the comment environment
\usepackage{verbatim}
\newenvironment{reference}{\comment}{\endcomment}
%\newenvironment{reference}{}{}
\newenvironment{slogan}{\comment}{\endcomment}
\newenvironment{history}{\comment}{\endcomment}

% For commutative diagrams we use Xy-pic
\usepackage[all]{xy}

% We use 2cell for 2-commutative diagrams.
\xyoption{2cell}
\UseAllTwocells

% We use multicol for the list of chapters between chapters
\usepackage{multicol}

% This is generally recommended for better output
\usepackage{lmodern}
\usepackage[T1]{fontenc}

% For cross-file-references

% Package for hypertext links:
\usepackage{hyperref}

% For any local file, say "hello.tex" you want to link to please

% Theorem environments.
%
\theoremstyle{plain}
\newtheorem{theorem}[subsection]{Theorem}
\newtheorem{proposition}[subsection]{Proposition}
\newtheorem{lemma}[subsection]{Lemma}

\theoremstyle{definition}
\newtheorem{definition}[subsection]{Definition}
\newtheorem{example}[subsection]{Example}
\newtheorem{exercise}[subsection]{Exercise}
\newtheorem{situation}[subsection]{Situation}

\theoremstyle{remark}
\newtheorem{remark}[subsection]{Remark}
\newtheorem{remarks}[subsection]{Remarks}

\numberwithin{equation}{subsection}

% Macros
%
\def\lim{\mathop{\mathrm{lim}}\nolimits}
\def\colim{\mathop{\mathrm{colim}}\nolimits}
\def\Spec{\mathop{\mathrm{Spec}}}
\def\Hom{\mathop{\mathrm{Hom}}\nolimits}
\def\Ext{\mathop{\mathrm{Ext}}\nolimits}
\def\SheafHom{\mathop{\mathcal{H}\!\mathit{om}}\nolimits}
\def\SheafExt{\mathop{\mathcal{E}\!\mathit{xt}}\nolimits}
\def\Sch{\mathit{Sch}}
\def\Mor{\mathop{\mathrm{Mor}}\nolimits}
\def\Ob{\mathop{\mathrm{Ob}}\nolimits}
\def\Sh{\mathop{\mathit{Sh}}\nolimits}
\def\NL{\mathop{N\!L}\nolimits}
\def\CH{\mathop{\mathrm{CH}}\nolimits}
\def\proetale{{pro\text{-}\acute{e}tale}}
\def\etale{{\acute{e}tale}}
\def\QCoh{\mathit{QCoh}}
\def\Ker{\mathop{\mathrm{Ker}}}
\def\Im{\mathop{\mathrm{Im}}}
\def\Coker{\mathop{\mathrm{Coker}}}
\def\Coim{\mathop{\mathrm{Coim}}}

% Boxtimes
%
\DeclareMathSymbol{\boxtimes}{\mathbin}{AMSa}{"02}

%
% Macros for moduli stacks/spaces
%
\def\QCohstack{\mathcal{QC}\!\mathit{oh}}
\def\Cohstack{\mathcal{C}\!\mathit{oh}}
\def\Spacesstack{\mathcal{S}\!\mathit{paces}}
\def\Quotfunctor{\mathrm{Quot}}
\def\Hilbfunctor{\mathrm{Hilb}}
\def\Curvesstack{\mathcal{C}\!\mathit{urves}}
\def\Polarizedstack{\mathcal{P}\!\mathit{olarized}}
\def\Complexesstack{\mathcal{C}\!\mathit{omplexes}}
% \Pic is the operator that assigns to X its picard group, usage \Pic(X)
% \Picardstack_{X/B} denotes the Picard stack of X over B
% \Picardfunctor_{X/B} denotes the Picard functor of X over B
\def\Pic{\mathop{\mathrm{Pic}}\nolimits}
\def\Picardstack{\mathcal{P}\!\mathit{ic}}
\def\Picardfunctor{\mathrm{Pic}}
\def\Deformationcategory{\mathcal{D}\!\mathit{ef}}

\setlength{\emergencystretch}{3em}
\begin{document}
\title[Stacks Project, Tag 058K]{1274 --- Pure short exact sequences over arbitrary rings are directed colimits of split sequences with finitely presented quotients}
\maketitle
\noindent Copyright (C) 2005--2025 Johan de Jong. Stacks Project authors. Source: \href{https://stacks.math.columbia.edu/tag/058K}{Tag 058K}.

\noindent Editorial difficulty note (not author text). Difficulty H1: The six-equivalence theorem connects tensor purity, solvability of finite relation systems, lifting from finitely presented modules and an explicit directed split-sequence presentation. Its finite-presentation approximation and compatible fiber-product construction exceed ordinary qualification exact-sequence manipulations..

\noindent Editorial provenance note (not author text). History: excerpt from revision \texttt{a0444\allowbreak 6e57e\allowbreak c1fbc\allowbreak 252a8\allowbreak 71afc\allowbreak ec775\allowbreak 2fb28\allowbreak 07b14}, algebra.tex, lines 19784--19947. Reference presentation was adapted; original-author context and core support blocks were retained or added. The mathematical wording is unchanged. Licensed under GNU FDL 1.2 or later; no invariant sections or cover texts. See ../licenses/COPYING. Context blocks reproduce author definitions and supporting results; they are not additional counted problems.

\begin{definition}
\label{definition-universally-injective}
Let $f: M \to N$ be a map of $R$-modules.  Then $f$ is called
{\it universally injective} if for every $R$-module $Q$, the map $f
\otimes_R \text{id}_Q: M \otimes_R Q \to N \otimes_R Q$
is injective.  A sequence $0 \to M_1 \to M_2 \to M_3
\to 0$ of $R$-modules is called {\it universally exact} if it is exact
and $M_1 \to M_2$ is universally injective.
\end{definition}

\begin{example}
\label{example-universally-exact}
Examples of universally exact sequences.
\begin{enumerate}
\item A split short exact sequence is universally exact since tensoring
commutes with taking direct sums.
\item The colimit of a directed system of universally exact sequences is
universally exact.  This follows from the fact that taking directed colimits is
exact and that tensoring commutes with taking colimits.  In particular the
colimit of a directed system of split exact sequences is universally exact.  We
will see below that, conversely, any universally exact sequence arises in this
way.
\end{enumerate}
\end{example}

\begin{lemma}
\label{lemma-module-colimit-fp}
Let $R$ be a ring and let $M$ be an $R$-module.
Then $M$ is the colimit of a directed system
$(M_i, \mu_{ij})$ of $R$-modules
with all $M_i$ finitely presented $R$-modules.
\end{lemma}

\begin{proof}
Consider any finite subset $S \subset M$ and any finite
collection of relations $E$ among the elements
of $S$. So each $s \in S$ corresponds to $x_s \in M$ and
each $e \in E$ consists of a vector
of elements $f_{e, s} \in R$ such that $\sum f_{e, s} x_s = 0$.
Let $M_{S, E}$ be the cokernel of the map
$$
R^{\# E} \longrightarrow R^{\# S}, \quad
(g_e)_{e\in E} \longmapsto (\sum g_e f_{e, s})_{s\in S}.
$$
There are canonical maps $M_{S, E} \to M$.
If $S \subset S'$ and if the elements of
$E$ correspond, via this map, to relations
in $E'$, then there is an obvious map
$M_{S, E} \to M_{S', E'}$ commuting with the
maps to $M$. Let $I$ be the set of pairs
$(S, E)$ with ordering by inclusion as above.
It is clear that the colimit of this directed system is $M$.
\end{proof}


\begin{theorem}
\label{theorem-universally-exact-criteria}
Let
$$
0 \to M_1 \xrightarrow{f_1} M_2 \xrightarrow{f_2} M_3 \to 0
$$
be an exact sequence of $R$-modules.  The following are equivalent:
\begin{enumerate}
\item The sequence $0 \to M_1 \to M_2 \to M_3
\to 0$ is universally exact.
\item For every finitely presented $R$-module $Q$, the sequence
$$
0 \to M_1 \otimes_R Q \to M_2 \otimes_R Q \to
M_3 \otimes_R Q \to 0
$$
is exact.
\item Given elements $x_i \in M_1$ $(i = 1, \ldots, n)$, $y_j \in M_2$ $(j = 1,
\ldots, m)$, and $a_{ij} \in R$ $(i = 1, \ldots, n, j = 1, \ldots, m)$ such that
for all $i$
$$
f_1(x_i) = \sum\nolimits_j a_{ij} y_j,
$$
there exists $z_j \in M_1$ $(j =1, \ldots, m)$ such that for all $i$,
$$
x_i = \sum\nolimits_j a_{ij} z_j .
$$
\item Given a commutative diagram of $R$-module maps
$$
\xymatrix{
R^n \ar[r] \ar[d] &  R^m \ar[d] \\
M_1 \ar[r]^{f_1}        &  M_2
}
$$
where $m$ and $n$ are integers, there exists a map $R^m \to M_1$ making
the top triangle commute.
\item For every finitely presented $R$-module $P$, the $R$-module
map $\Hom_R(P, M_2) \to \Hom_R(P, M_3)$ is surjective.
\item The sequence $0 \to M_1 \to M_2 \to M_3
\to 0$ is the colimit of a directed system of split exact sequences of
the form
$$
0 \to M_{1} \to M_{2, i} \to M_{3, i} \to 0
$$
where the $M_{3, i}$ are finitely presented.
\end{enumerate}
\end{theorem}

\begin{proof}
Obviously (1) implies (2).

\medskip\noindent
Next we show (2) implies (3).  Let $f_1(x_i) = \sum_j a_{ij} y_j$ be relations
as in (3).  Let $(d_j)$ be a basis for $R^m$, $(e_i)$ a basis for $R^n$, and
$R^m \to R^n$ the map given by $d_j \mapsto \sum_i a_{ij} e_i$.
Let $Q$ be the cokernel of $R^m \to R^n$.  Then tensoring
$R^m \to R^n \to Q \to 0$ by the map $f_1: M_1 \to M_2$, we get a
commutative diagram
$$
\xymatrix{
M_1^{\oplus m} \ar[r] \ar[d] & M_1^{\oplus n} \ar[r] \ar[d] & M_1 \otimes_R Q
\ar[r] \ar[d] & 0 \\
M_2^{\oplus m} \ar[r] & M_2^{\oplus n} \ar[r] & M_2 \otimes_R Q \ar[r] & 0
}
$$
where $M_1^{\oplus m} \to M_1^{\oplus n}$ is given by
$$
(z_1, \ldots, z_m) \mapsto
(\sum\nolimits_j a_{1j} z_j, \ldots, \sum\nolimits_j a_{nj} z_j),
$$
and $M_2^{\oplus m} \to M_2^{\oplus n}$ is given similarly.  We want to
show $x = (x_1, \ldots, x_n) \in M_1^{\oplus n}$ is in the image of $M_1^{\oplus
m} \to M_1^{\oplus n}$.  By (2) the map $M_1 \otimes Q \to M_2
\otimes Q$ is injective, hence by exactness of the top row it is enough to show
$x$ maps to $0$ in $M_2 \otimes Q$, and so by exactness of the bottom row it is
enough to show the image of $x$ in $M_2^{\oplus n}$ is in the image of
$M_2^{\oplus m} \to M_2^{\oplus n}$.  This is true by assumption.

\medskip\noindent
Condition (4) is just a translation of (3) into diagram form.

\medskip\noindent
Next we show (4) implies (5). Let $\varphi : P \to M_3$ be a map from a
finitely presented $R$-module $P$.  We must show that $\varphi$ lifts to a map
$P \to M_2$.  Choose a presentation of $P$,
$$
R^n \xrightarrow{g_1} R^m \xrightarrow{g_2} P \to 0.
$$
Using freeness of $R^n$ and $R^m$, we can construct $h_2: R^m \to M_2$
and then $h_1: R^n \to M_1$ such that the following diagram commutes
$$
\xymatrix{
         & R^n \ar[r]^{g_1} \ar[d]^{h_1} & R^m \ar[r]^{g_2} \ar[d]^{h_2} & P
\ar[r] \ar[d]^{\varphi} & 0 \\
0 \ar[r] & M_1 \ar[r]^{f_1} & M_2 \ar[r]^{f_2} & M_3 \ar[r] & 0 .
}
$$
By (4) there is a map $k_1: R^m \to M_1$ such that $k_1 \circ g_1 =
h_1$.  Now define $h'_2: R^m \to M_2$ by $h_2' = h_2 - f_1 \circ k_1$.
Then
$$
h'_2 \circ g_1 = h_2 \circ g_1 - f_1 \circ k_1 \circ g_1 =
h_2 \circ g_1 - f_1 \circ h_1 = 0 .
$$
Hence by passing to the quotient $h'_2$ defines a map $\varphi': P \to
M_2$ such that $\varphi' \circ g_2 = h_2'$.  In a diagram, we have
$$
\xymatrix{
R^m \ar[r]^{g_2} \ar[d]_{h'_2} & P \ar[d]^{\varphi} \ar[dl]_{\varphi'} \\
M_2 \ar[r]^{f_2} & M_3.
}
$$
where the top triangle commutes.  We claim that $\varphi'$ is the desired lift,
i.e.\ that $f_2 \circ \varphi' = \varphi$.  From the definitions we have
$$
f_2 \circ \varphi' \circ g_2 = f_2 \circ h'_2 =
f_2 \circ h_2 - f_2 \circ f_1 \circ k_1 = f_2 \circ h_2 =
\varphi \circ g_2.
$$
Since $g_2$ is surjective, this finishes the proof.

\medskip\noindent
Now we show (5) implies (6). Write $M_{3}$ as the colimit of a directed system
of finitely presented modules $M_{3, i}$, see
Lemma \ref{lemma-module-colimit-fp}. Let $M_{2, i}$ be the fiber product
of $M_{3, i}$ and $M_{2}$ over $M_{3}$---by definition this is the submodule of
$M_2 \times M_{3, i}$ consisting of elements whose two projections onto $M_3$
are equal. Let $M_{1, i}$ be the kernel of the projection
$M_{2, i} \to M_{3, i}$. Then we have a directed system of exact sequences
$$
0 \to M_{1, i} \to M_{2, i} \to M_{3, i} \to 0,
$$
and for each $i$ a map of exact sequences
$$
\xymatrix{
0  \ar[r] & M_{1, i} \ar[d] \ar[r]  &  M_{2, i}  \ar[r] \ar[d] & M_{3, i} \ar[d]
\ar[r] & 0 \\
0  \ar[r] & M_{1} \ar[r]  &  M_{2}  \ar[r] & M_{3} \ar[r] & 0
}
$$
compatible with the directed system.  From the definition of the fiber product
$M_{2, i}$, it follows that the map $M_{1, i} \to M_1$ is an isomorphism.
 By (5) there is a map $M_{3, i} \to M_{2}$ lifting $M_{3, i} \to
M_3$, and by the universal property of the fiber product this gives rise to a
section of $M_{2, i} \to M_{3, i}$.  Hence the sequences
$$
0 \to M_{1, i} \to M_{2, i} \to M_{3, i} \to 0
$$
split.  Passing to the colimit, we have a commutative diagram
$$
\xymatrix{
0  \ar[r] & \colim M_{1, i} \ar[d]^{\cong} \ar[r]  &  \colim M_{2, i}
 \ar[r]
\ar[d] & \colim M_{3, i} \ar[d]^{\cong} \ar[r] & 0 \\
0  \ar[r] & M_{1} \ar[r]  &  M_{2}  \ar[r] & M_{3} \ar[r] & 0
}
$$
with exact rows and outer vertical maps isomorphisms.  Hence $\colim
M_{2, i}
\to M_2$ is also an isomorphism and (6) holds.

\medskip\noindent
Condition (6) implies (1) by
Example \ref{example-universally-exact} (2).
\end{proof}
\end{document}
